%=============================================================
% Paper 1: Davis & Haltiwanger (1992), QJE
%=============================================================
\chapter{Davis \& Haltiwanger (1992)}

\begin{paperbox}{Gross Job Creation, Gross Job Destruction, and Employment Reallocation}
\textbf{Authors:} Steven J.\ Davis \& John Haltiwanger\\[3pt]
\textbf{Journal:} \textit{Quarterly Journal of Economics} 107(3): 819--863\\[3pt]
\textbf{Year:} 1992 \quad\textbf{Field:} Labour economics; Macroeconomics\\[3pt]
\textbf{Classification:} Empirical (descriptive measurement)
\end{paperbox}

%------------------------------------------------------------
\section{Research Question}

What is the magnitude and structure of gross job creation and destruction in
U.S.\ manufacturing, how does job reallocation relate to worker reallocation,
and what accounts for simultaneous job creation and destruction within
narrowly defined sectors of the economy?

%------------------------------------------------------------
\section{Motivation}

Prior to this paper, macroeconomics focused almost exclusively on \emph{net}
employment changes, treating aggregate employment as if it reflected
homogeneous expansions and contractions across the economy. The large gross
worker flows documented in the CPS literature (Clark--Summers, Abowd--Zellner)
remained unexplained in structural terms. It was unclear whether large worker
flows primarily reflected temporary layoffs and recalls across a fixed set of
jobs, or whether a large portion was driven by the genuine destruction and
creation of employment opportunities. The answer matters: if most worker
mobility is driven by job reallocation, the character, timing, and welfare cost
of unemployment are very different from a world where workers simply shuffle
across stable jobs.

%------------------------------------------------------------
\section{Main Contribution}

The paper makes four major contributions.

\begin{enumerate}
  \item \textbf{Measurement framework.} Constructs the DH gross job flow
    framework---the definitions of POS (job creation), NEG (job destruction),
    SUM (job reallocation), and the Davis--Haltiwanger symmetric growth rate
    measure---which became the field standard. U.S.\ manufacturing averaged
    9.2\% job creation and 11.3\% job destruction per year (1972--1986).

  \item \textbf{Job--worker reallocation link.} Demonstrates that
    35--56\% of all worker reallocation between employment states is directly
    driven by job reallocation: shifts in the distribution of employment
    opportunities across plants are a dominant source of worker mobility.

  \item \textbf{Persistence and concentration.} Job creation and destruction
    are largely \emph{persistent} (one-year rates: 68\% for creation, 81\%
    for destruction) and \emph{concentrated} (77\% of destruction from plants
    shrinking by $>$20\%), ruling out temporary layoff/recall as the dominant
    mechanism.

  \item \textbf{Cross-sectional heterogeneity.} Documents systematic
    variation in job reallocation by plant age, size, ownership type, and
    region; shows that between-sector shifts account for essentially none
    ($\approx$12\% at the four-digit level) of excess job reallocation.
\end{enumerate}

%------------------------------------------------------------
\section{Relationship to the Literature}

The paper sits at the intersection of three literatures: (i) gross worker
flows (Clark--Summers 1979; Hall 1982; Blanchard--Diamond 1990); (ii)
firm/plant-level employment dynamics (Jovanovic 1982; Evans 1987; Dunne,
Roberts, Samuelson 1989); and (iii) macroeconomic business-cycle theory
(Caballero 1990; Davis--Haltiwanger 1990). Among U.S.\ studies, only Leonard
(1987) and Dunne, Roberts, Samuelson (1989) had previously measured gross job
flows at the plant level, but with less frequency, shorter samples, and
narrower coverage. The paper directly motivates Hopenhayn (1992) and much of
the ensuing structural literature on firm dynamics.

%------------------------------------------------------------
\section{Theoretical Framework and Economic Mechanism}

The paper is primarily measurement-oriented. The conceptual apparatus rests on
three mechanisms proposed in the literature to explain simultaneous job
creation and destruction:

\begin{itemize}
  \item \textbf{Passive learning and selection} (Jovanovic 1982): plants
    learn their efficiency type; those discovering unfavourable types exit
    while survivors grow.
  \item \textbf{Active learning / vintage effects} (Ericson--Pakes 1989):
    stochastic investment outcomes generate ongoing idiosyncratic uncertainty,
    sustaining perpetual entry and exit.
  \item \textbf{Idiosyncratic cost/demand disturbances} (Hopenhayn 1989;
    Davis--Haltiwanger 1990): even homogeneous plants within sectors face
    local idiosyncratic shocks, driving simultaneous expansion and contraction.
\end{itemize}

\subsection*{Measurement Definitions}

The growth rate of establishment $e$ at time $t$ is defined as:
\begin{equation}
  g_{et} = \frac{E_{et} - E_{e,t-1}}{\tfrac{1}{2}(E_{et} + E_{e,t-1})},
  \label{eq:dh_growth}
\end{equation}
which lies in $[-2,\,2]$ and symmetrically handles births ($g=2$) and
deaths ($g=-2$). Gross job creation and destruction rates for sector $s$
at time $t$ are:
\begin{align}
  \text{POS}_{st} &= \sum_{\substack{e \in \mathcal{E}_{st}\\ g_{et}>0}}
    \frac{x_{et}}{X_{st}}\, g_{et}, \label{eq:pos} \\[4pt]
  \text{NEG}_{st} &= \sum_{\substack{e \in \mathcal{E}_{st}\\ g_{et}<0}}
    \frac{x_{et}}{X_{st}}\, |g_{et}|, \label{eq:neg}
\end{align}
where $x_{et}=\tfrac{1}{2}(E_{et}+E_{e,t-1})$ is the size metric and
$X_{st}$ is the sector-level analogue. The job reallocation rate is
$\text{SUM}_{st}=\text{POS}_{st}+\text{NEG}_{st}$. The required worker
reallocation is bounded by
$\text{MAX}_{st}=\max\{\text{POS}_{st},\text{NEG}_{st}\}$ (lower bound)
and $\text{SUM}_{st}$ (upper bound).

%------------------------------------------------------------
\section{Data}

\begin{itemize}
  \item \textbf{Dataset:} Longitudinal Research Datafile (LRD), U.S.\ Census
    Bureau.
  \item \textbf{Coverage:} All manufacturing establishments with 5+ employees
    sampled in the Annual Survey of Manufactures (ASM).
  \item \textbf{Sample:} $\approx$860,000 annual observations on
    $\approx$160,000 establishments; 11 years (1972--1986, excluding 1974,
    1979, 1984 due to ASM panel break-years).
  \item \textbf{Strengths:} Probability-based sample; births incorporated
    into ongoing panels; careful distinction between ownership transfers and
    actual births/deaths; certainty stratum covers $\approx$two-thirds of
    manufacturing employment ($\geq$250 employees).
  \item \textbf{Limitation:} Five-year panel structure makes non-certainty
    (small) establishments harder to track across panel boundaries; point-in-time
    (March) employment understates true intra-year flows.
\end{itemize}

%------------------------------------------------------------
\section{Empirical Strategy}

Pure descriptive measurement. The strategy involves:
\begin{enumerate}
  \item Computing annual POS, NEG, SUM, MAX at the establishment level and
    aggregating using employment shares as weights.
  \item Tabulating \emph{persistence}: what fraction of jobs created (destroyed)
    in year $t$ remain created (destroyed) one and two years later.
  \item Cross-tabulating gross flow rates by industry, size class, age group,
    ownership type, and geographic region.
  \item Bounding required worker reallocation (SUM upper bound; MAX lower
    bound) and comparing against CPS-based total worker reallocation.
\end{enumerate}

%------------------------------------------------------------
\section{Identification Strategy}

\textbf{This is a descriptive paper with no causal identification.} The
approach is measurement and decomposition. For the passive-learning exercise,
the authors use an indirect decomposition: plants older than some age $T^{*}$
are assumed to have fully resolved initial uncertainty, so the excess
reallocation of younger plants relative to this benchmark is attributed to
learning. This requires two identifying assumptions: (i) within-industry,
long-run job reallocation rates are stable across age cohorts for old plants;
(ii) excess reallocation of young plants is not driven by other factors
that diminish mechanically with age.

%------------------------------------------------------------
\section{Main Results}

\subsection*{Core Measurement (Table~I)}

\begin{center}
\begin{tabular}{lcc}
\toprule
\textbf{Statistic} & \textbf{Value} & \textbf{Period} \\
\midrule
Average annual job creation rate (POS) & 9.2\% & 1972--1986 \\
Average annual job destruction rate (NEG) & 11.3\% & 1972--1986 \\
Average job reallocation rate (SUM) & 20.5\% & 1972--1986 \\
Range of job reallocation rate & 17.3--23.3\% & \\
$\rho(\text{POS},\,\text{NEG})$ & $-0.864$ & \\
$\rho(\text{NET},\,\text{SUM})$ & $-0.565$ & \\
\bottomrule
\end{tabular}
\end{center}

\subsection*{Cross-industry variation (Table~II)}

Job reallocation ranges from 14.0\% (Tobacco) to 28.8\% (Lumber) across
2-digit industries. Every industry exhibits simultaneous creation and
destruction; excess reallocation ranges from 9.8\% to 20.6\%.

\subsection*{Persistence (Table~III)}

\begin{center}
\begin{tabular}{lcc}
\toprule
 & \textbf{Job Creation} & \textbf{Job Destruction} \\
\midrule
One-year persistence (mean) & 68\% & 81\% \\
Two-year persistence (mean) & 50\% & 73\% \\
\bottomrule
\end{tabular}
\end{center}

\subsection*{Plant characteristics (Table~IV)}

\begin{itemize}
  \item \textbf{Age:} reallocation falls from 48\% at age 1 to 16\% at
    15+ years---the sharpest gradient in the data.
  \item \textbf{Size:} reallocation falls from 30\% (1--99 employees) to
    14\% (1000+ employees).
  \item \textbf{Ownership:} single-unit plants 28\%; multi-unit 18\%.
\end{itemize}

\subsection*{Job--worker reallocation link}

Total worker reallocation $\approx$36.8\% of employment annually (CPS).
Job reallocation accounts for 35--56\% of total worker reallocation.

%------------------------------------------------------------
\section{Economic Magnitude and Interpretation}

The magnitude is striking. In an average year, employment positions equivalent
to one-fifth of the manufacturing workforce are either newly created or newly
destroyed. Even in boom years, gross job destruction exceeds 6\%. The
persistence results imply that most job creation and destruction represents
real, durable changes---not temporary layoff/recall cycles. The concentration
result (77\% of destruction from plants shrinking $>$20\%) means most job
losses cannot be absorbed by normal attrition and require actual layoffs.
These findings establish that standard representative-agent macro models,
which focus entirely on net employment, ignore enormous amounts of
economically significant labour market activity.

%------------------------------------------------------------
\section{Mechanisms}

The paper evaluates four potential explanations:

\begin{enumerate}
  \item \textbf{Between-sector employment shifts:} Even at the 4-digit
    industry level (450 industries), between-sector shifts account for only
    $\approx$12\% of excess job reallocation. Sectoral-shock explanations
    are definitively ruled out.

  \item \textbf{Passive learning:} Accounts for only 11--13\% of \emph{total}
    reallocation, but explains one-third to one-half of the \emph{variation}
    in reallocation rates across age, size, and ownership groups.

  \item \textbf{Plant-level idiosyncratic disturbances:} By elimination, the
    dominant mechanism must be ongoing plant-level idiosyncratic shocks
    orthogonal across plants even within the same 4-digit industry.

  \item \textbf{Cyclical variation:} The countercyclical pattern reflects
    time variation in the \emph{magnitude} of idiosyncratic shocks, not
    compositional shifts. Patterns are most pronounced in older, larger,
    multi-unit plants.
\end{enumerate}

%------------------------------------------------------------
\section{Robustness Checks}

\begin{itemize}
  \item Consistency of gross flow estimates with LRD-period BLS Business
    Employment Dynamics data.
  \item Consistency of CPS-based gross worker flow adjustments using both
    Abowd--Zellner and Poterba--Summers error matrices.
  \item Results hold across multiple base years for persistence calculations.
  \item Cross-industry variation results replicated across individual years.
\end{itemize}

\textbf{Key limitation acknowledged:} The LRD excludes establishments with
$<$5 employees. Leonard (1987) finds nonmanufacturing reallocation is
$\approx$28\% higher than manufacturing, suggesting economy-wide reallocation
substantially exceeds the manufacturing estimates.

%------------------------------------------------------------
\section{Limitations}

\begin{enumerate}
  \item \textbf{Sector coverage:} Confined to U.S.\ manufacturing (declining
    over the period). Service sectors, with different dynamics, are excluded.
  \item \textbf{Annual frequency:} March-to-March data miss intra-year job
    flows; reported annual flows understate true quarterly or monthly flows.
  \item \textbf{Point-in-time employment:} Measures lower bounds on true
    gross flows.
  \item \textbf{No causal identification:} Documents regularities but cannot
    explain \emph{why} reallocation is countercyclical.
  \item \textbf{Learning decomposition is fragile:} Requires identifying
    assumptions about when reallocation stabilises at its long-run level.
\end{enumerate}

%------------------------------------------------------------
\section{Policy Implications}

\begin{itemize}
  \item High reallocation rates imply unemployment is typically a structural
    transition phenomenon; job-search assistance and training are more relevant
    than aggregate demand stimulus.
  \item Countercyclical reallocation supports ``cleansing recession'' arguments
    against bailouts that preserve misallocated resources.
  \item Firing-cost regulations would reduce both creation and reallocation,
    potentially lowering long-run productivity.
  \item Large role of young and small plants in total reallocation suggests
    that entry costs (regulations, administrative burdens) have first-order
    effects on allocative efficiency.
\end{itemize}

%------------------------------------------------------------
\section{Overall Assessment}

This is a landmark empirical paper that created a field. The DH gross job flow
framework is now used worldwide; the LRD-based estimates have been replicated
and extended to dozens of countries and sectors. The paper's weaknesses are
largely a function of when it was written: limited coverage of small plants,
manufacturing-only scope, and no causal identification. The paper does not
attempt to distinguish whether countercyclical reallocation is efficient
(``cleansing'') or inefficient (``sullying''), a distinction that would motivate
a decade of subsequent theoretical work.

\textbf{Quality: Exceptional. Required reading for anyone working on labour
dynamics, firm dynamics, or macroeconomic fluctuations.}

%------------------------------------------------------------
\section{Relevance to My Research}

Directly foundational for RDC-based research on job flows in Canada. The
methodology (plant-level longitudinal data, decomposition of gross flows by
plant characteristics) is exactly the approach applicable to LEAP or LEED data.
The key measurements would be directly replicable with Canadian data and would
reveal whether the magnitudes found in U.S.\ manufacturing generalise to the
Canadian context---including resource-sector-heavy Western provinces.

%------------------------------------------------------------
\section{Potential Research Extensions}

\begin{itemize}
  \item Replicate the DH framework for Canadian manufacturing and services
    using LEAP/LEED; test whether Canadian job reallocation is higher or lower
    than U.S., particularly around resource price cycles.
  \item Compare job reallocation dynamics before and after major trade
    liberalisations (FTA 1989, NAFTA 1994, CETA).
  \item Examine whether countercyclical job reallocation intensified during
    the 2008--09 recession and COVID-19 in Canada.
\end{itemize}

%------------------------------------------------------------
\section{Research Gaps}

\begin{itemize}
  \item \textbf{Causal mechanism:} Does countercyclical reallocation represent
    efficient cleansing or inefficient disruption? The paper is silent on this
    welfare question.
  \item \textbf{Worker outcomes:} Job destruction rates tell us about
    positions, not about workers. Earnings losses, nonemployment duration, and
    occupational switching are entirely outside the scope.
  \item \textbf{Service sector:} U.S.\ nonmanufacturing dynamics at the plant
    level remain underexplored.
\end{itemize}

%------------------------------------------------------------
\section{Methodological Lessons}

\begin{enumerate}
  \item The DH growth rate measure $g_{et}=\Delta E/\text{average}(E)$ is
    superior to the log-difference because it handles births and deaths
    symmetrically within a unified framework.
  \item Persistence measures (fraction of newly created/destroyed jobs still
    present one or two years later) are essential for distinguishing transitory
    from permanent employment changes.
  \item Bounding worker reallocation from above (SUM) and below (MAX) is a
    simple and powerful diagnostic for the connection between job and worker
    flows.
  \item Weighting establishment-level growth rates by employment share
    (rather than equally) is critical for ensuring aggregate implications
    reflect the behaviour of large employers.
\end{enumerate}

%------------------------------------------------------------
\section*{Five-Sentence Summary}

Davis and Haltiwanger (1992) use the U.S.\ Census Bureau's Longitudinal
Research Datafile to document, for the first time comprehensively, that gross
job creation and destruction in U.S.\ manufacturing averaged 9.2\% and 11.3\%
of employment per year over 1972--1986---flows an order of magnitude larger
than net employment changes. They show that job reallocation is persistent
(68--81\% one-year persistence), concentrated (77\% of destruction from plants
shrinking $>$20\%), and accounts for 35--56\% of all worker reallocation,
establishing that shifts in the distribution of employment opportunities are a
primary driver of worker mobility. Cross-sectional analysis reveals that job
reallocation rates decline sharply with plant age (48\% at age~1; 16\% at
15+~years) and size, and that virtually none of excess reallocation can be
attributed to between-sector employment shifts even at the 4-digit industry
level. The countercyclical pattern of job reallocation is driven by time
variation in the magnitude of idiosyncratic plant-level shocks, primarily at
older and larger establishments. The paper establishes the conceptual and
measurement framework---the DH gross job flow decomposition---that has become
the foundation for all subsequent empirical and theoretical work on firm
dynamics, labour market flows, and macroeconomic reallocation.

\vspace{6pt}
\begin{quickref}
\begin{tabularx}{\linewidth}{@{}lX@{}}
\textbf{Key contribution} &
  First comprehensive plant-level measurement of gross job flows; establishes
  that job reallocation is large, persistent, and drives a substantial fraction
  of worker mobility. \\[4pt]
\textbf{Key weakness} &
  Descriptive only; confined to manufacturing; annual frequency understates
  true flow magnitudes; no causal identification of the countercyclical
  reallocation mechanism. \\[4pt]
\textbf{Most interesting gap} &
  What happens to workers displaced by job destruction---the paper is silent
  on earnings losses, nonemployment duration, and occupational switching. \\[4pt]
\textbf{Publishable extension} &
  Replicate the DH framework for Canadian data (LEAP/LEED) across the full
  economy, with particular attention to resource-sector dynamics and commodity
  price cycles on job reallocation patterns.
\end{tabularx}
\end{quickref}
